\documentclass[lualatex,book,paper=a4paper,jafontsize=11.5pt,jafontscale=1.05]{jlreq}

%フォント:
\usepackage{luatexja-fontspec}
\usepackage{hack-fonts}

%表紙用:
\usepackage{hack-cover}

% リンク
\usepackage[hidelinks]{hyperref}

% 画像
\usepackage{graphicx}


\title{ハッカソン}
\subtitle{～計画書と設計書～}
\team{チーム○○}
\author{
	2xG1xxx : FFFF NNNN\\
        2xG1xxx : FFFF NNNN\\
	2xG1xxx : FFFF NNNN\\
        2xG1xxx : FFFF NNNN\\
        2xG1xxx : FFFF NNNN\\
        2xG1xxx : FFFF NNNN 
        }
\date{\today}
\version{x.x}

%%%%%%%%%%%%%%%%%%%%%%%%%%%%%%%%%%%%%%%%%%%%%%%%
\begin{document}
\maketitle

\setcounter{tocdepth}{2}
\tableofcontents \thispagestyle{empty}

%%%%%%%%%%%%%%%%%%%%%%%%
%%%%%%%%%%%%%%%%%%%%%%%%
%%%%%%%%%%%%%%%%%%%%%%%%
\chapter{プロジェクトの計画書} \label{chap:project-plan}\setcounter{page}{1}


%%%%%%%%%%%%%%%%%%%%%%%%
\section{プロジェクトの概要}
%%%%%%%%%%%%%%%%%%%%%%%%
\subsection{プロジェクトの題目}

%%%%%%%%%%%%%%%%%%%%%%%%
\subsection{プロジェクトの目的}
プロジェクトの目的と最終的なゴール，解決すべき課題を記述する．

%%%%%%%%%%%%%%%%%%%%%%%%
\subsection{既存の技術とプロジェクトの独自性}

%%%%%%%%%%%%%%%%%%%%%%%%
%%%%%%%%%%%%%%%%%%%%%%%%
\section{スコープ・マネジメント}
%%%%%%%%%%%%%%%%%%%%%%%%
\subsection{成果物}

%%%%%%%%%%%%%%%%%%%%%%%%
\subsection{プロジェクトの対象範囲と非対象範囲}

%%%%%%%%%%%%%%%%%%%%%%%%
\subsection{機能分解構造FBS}

%%%%%%%%%%%%%%%%%%%%%%%%
\subsection{製品分解構造PBS}

%%%%%%%%%%%%%%%%%%%%%%%%
\subsection{FBSとPBSの対応関係}


%%%%%%%%%%%%%%%%%%%%%%%%
%%%%%%%%%%%%%%%%%%%%%%%%
\section{作業計画}
%%%%%%%%%%%%%%%%%%%%%%%%
\subsection{作業分解構造WBSと管理番号}

%%%%%%%%%%%%%%%%%%%%%%%%
\subsection{アローダイアグラムとクリティカルパス}

%%%%%%%%%%%%%%%%%%%%%%%%
\subsection{役割分担}

%%%%%%%%%%%%%%%%%%%%%%%%
\subsection{ガントチャート}


%%%%%%%%%%%%%%%%%%%%%%%%
%%%%%%%%%%%%%%%%%%%%%%%%
\section{検証と妥当性確認: V\&V計画}
%%%%%%%%%%%%%%%%%%%%%%%%
\subsection{プロジェクトの成果物に対する有効指標MOE}

%%%%%%%%%%%%%%%%%%%%%%%%
\subsection{有効指標MOEに対応する性能指標MOP}

%%%%%%%%%%%%%%%%%%%%%%%%
\subsection{システムテストで評価する性能指標MOP}

%%%%%%%%%%%%%%%%%%%%%%%%
\subsection{結合テストで評価する性能指標MOP}

%%%%%%%%%%%%%%%%%%%%%%%%
\subsection{単体テストで評価する性能指標MOP}

%%%%%%%%%%%%%%%%%%%%%%%%
\subsection{プロジェクト中に監視する技術性能指標TPM}

%%%%%%%%%%%%%%%%%%%%%%%%
\subsection{プロジェクト中に監視する指標TPM}


%%%%%%%%%%%%%%%%%%%%%%%%
%%%%%%%%%%%%%%%%%%%%%%%%
\section{資源・調達管理}


%%%%%%%%%%%%%%%%%%%%%%%%
%%%%%%%%%%%%%%%%%%%%%%%%
\section{リスク管理}


%%%%%%%%%%%%%%%%%%%%%%%%
%%%%%%%%%%%%%%%%%%%%%%%%
%%%%%%%%%%%%%%%%%%%%%%%%
\chapter{システムの基本設計}

設計書は作成する成果物によって変化します．
適宜修正して利用してください．


%%%%%%%%%%%%%%%%%%%%%%%%
\section{機能一覧}
機能分解構造FBSの要素の説明

%%%%%%%%%%%%%%%%%%%%%%%%
\section{システム構成図}
ハードウェア，ソフトウェアなどの構成図を記載する．

%%%%%%%%%%%%%%%%%%%%%%%%
\section{操作インターフェース設計}
画面やスイッチなどのインターフェースの設計図を記載する

%%%%%%%%%%%%%%%%%%%%%%%%
\section{状態遷移図}


%%%%%%%%%%%%%%%%%%%%%%%%
%%%%%%%%%%%%%%%%%%%%%%%%
%%%%%%%%%%%%%%%%%%%%%%%%
\chapter{システムの詳細設計}

設計書は作成する成果物によって変化します．
適宜修正して利用してください．


%%%%%%%%%%%%%%%%%%%%%%%%
\section{部品一覧}
製品分解構造PBSの要素の説明

%%%%%%%%%%%%%%%%%%%%%%%%
%%%%%%%%%%%%%%%%%%%%%%%%
\section{ハードウェア設計}

%%%%%%%%%%%%%%%%%%%%%%%%
%%%%%%%%%%%%%%%%%%%%%%%%
\section{ソフトウェア設計}

%%%%%%%%%%%%%%%%%%%%%%%%
\subsection{ファイル構成}

%%%%%%%%%%%%%%%%%%%%%%%%
\subsection{オブジェクト設計(クラス図)}

%%%%%%%%%%%%%%%%%%%%%%%%
\subsection{関数設計}


%%%%%%%%%%%%%%%%%%%%%%%%
%%%%%%%%%%%%%%%%%%%%%%%%
\section{処理フローの設計}


%%%%%%%%%%%%%%%%%%%%%%%%
%%%%%%%%%%%%%%%%%%%%%%%%
\section{テストの設計}
どのようにテストを行うのか，例えば，テストデータの設計などを記述．


\end{document}